On the domain $\Omega=\{(x,y)\in\mathbb{R}^2 \mid y>0\}$
the linear partial differential operator
\begin{equation}
L=
\frac{\partial^2}{\partial x^2}
-
4y
\frac{\partial^2}{\partial y^2}
\label{90000034:operator}
\end{equation}
of second order is given.
\begin{teilaufgaben}
\item
Is the operator $L$ elliptic, parabolic or hyperbolic?
\item
Find the part of the boundary, where boundary values can influence the
value of the solution in the point $P=(1,1)$.
In other words: assuming the the $y$-direction is a timelike 
coordinate, which part of the boundary is in the past of point $P$?
\end{teilaufgaben}

\begin{loesung}
\begin{figure}[h]
\definecolor{darkred}{rgb}{0.8,0,0}
\definecolor{darkgreen}{rgb}{0,0.6,0}
\centering
\begin{tikzpicture}[>=latex,thick,scale=4]
\draw[->] (0,-0.1) -- (0,1.5) coordinate[label={$y$}];
\draw[->] (-0.5,0) -- (2.5,0) coordinate[label={$x$}];
\fill[color=blue!20]
	plot[domain=0:1,samples=20] ({2-\x},{\x^2})
	--
	plot[domain=0:1,samples=20] ({1-\x},{(1-\x)^2})
	--
	cycle;
\begin{scope}
\clip (-0.45,0) rectangle (2.45,1.45);
\foreach \c in {-3,-2.8,...,2.4}{
	\draw[color=darkred] plot[domain=0:1.3,samples=20] ({\x+\c},{\x^2});
}
\foreach \c in {-0.4,-0.2,...,3.6}{
	\draw[color=darkgreen] plot[domain=1.3:0,samples=20] ({-\x+\c},{\x^2});
}
\end{scope}
\fill[color=blue] (1,1) circle[radius=0.02];
\draw[color=orange,line width=1.2pt] (0,0) -- (2,0);
\draw (1,-0.02) -- (1,0.02);
\node at (1,0) [below] {$1$};
\draw (2,-0.02) -- (2,0.02);
\node at (2,0) [below] {$2$};
\node[color=blue] at (1,1) [right] {$P$};
\draw (-0.02,1) -- (0.02,1);
\node at (0,1) [left] {$1$};
\end{tikzpicture}
\caption{Characteristics for the differential operator
\eqref{90000034:operator}
\label{90000034:fig}}
\end{figure}
\begin{teilaufgaben}
\item
The symbol matrix is
\[
A
=
\begin{pmatrix}
1&0\\
0&-4y
\end{pmatrix},
\]
which has the determinant
\[
\det A
=
-4y < 0
\]
as $y>0$ on $\Omega$,
so the operator definitely is hyperbolic.
\item
The equation for the characteristics can be read off of the symbol matrix,
it is
\[
\dot{y}^2 -4y\dot{x}^2 = 0.
\]
Dividing by $\dot{x}$ und using $y'=\dot{y}/\dot{x}$ gives
\[
y^{\prime 2} -4y =0.
\]
This differential equation has the factorization
\[
y^{\prime 2} - 4y
=
(y'-2\!\sqrt{y})(y'+2\!\sqrt{y})
=
0,
\]
which leads to the differential equations
\begin{align*}
y'
&=
2
\!\sqrt{y}
&&\text{and}&
y'
&=
-
2
\!\sqrt{y}
\intertext{with solutions}
\frac{y'}{2\!\sqrt{y}} &= 1
&&&
\frac{y'}{2\!\sqrt{y}} &= -1
\\
\int\frac{dy}{2\!\sqrt{y}} &= \int \,dx
&&&
\int\frac{dy}{2\!\sqrt{y}} &= -\int \,dx
\\
\sqrt{y} &=  x + C
&&&
\sqrt{y} &= -x + D
\\
y &= (x+C)^2
&&&
y &= (-x+D)^2.
\end{align*}
Thus the characteristics are parabolas that touch the $x$-axis, drawn
in figure~\ref{90000034:fig}.
In particular, the characteristics through the point $P=(1,1)$
correspond to the paramters $C=1$ and $D=1$, i.~e.~they start from
the boundary points $(0,0)$ and $(2,0)$.
Thus the boundary values in the points in the interval
$\{ (x,0)\mid 0\le x\le 2\}$ of the $x$-axis, shown in orange in
figure~\ref{90000034:fig}, can influence
the value of the solution in the point $P=(1,1)$.
\qedhere
\end{teilaufgaben}
\end{loesung}


\begin{bewertung}
Symbol matrix ({\bf A}) 1 point,
operator is hyperbolic ({\bf H}) 1 point,
differential equation of characteristics ({\bf D}) 1 point,
factorization or other method to convert them into two first order
differential equations ({\bf F}) 1 point,
characteristic curves ({\bf C}) 1 point,
domain of influence for the point $P$ ({\bf P}) 1 point.
\end{bewertung}

